\documentclass[11pt,a4paper]{article}
\usepackage[margin=25mm]{geometry}
\usepackage{amsmath,amssymb,amsthm,booktabs,array}
\newtheorem{theorem}{Theorem}
\newtheorem{lemma}[theorem]{Lemma}
\newcommand{\Sz}{\operatorname{Sz}}
\newcommand{\Tr}{\operatorname{Tr}}
\newcommand{\F}{\mathbb F}
\title{Descent of a Suzuki face through an odd field extension}
\author{Human Compression project: bounded F003 continuation}
\date{}
\begin{document}\maketitle

We complete the case in which one genuine face is $\Sz(8)$ and the
field of the other face contains $\F_8$. The two previous small-face
notes, including their full written proofs and exact certificate
interfaces, remain premises. Their classical Suzuki structure and
earlier reductions remain attributed to the sources recorded there.
The additional argument is a field-descent criterion for two free
matrix entries. It uses an affine consequence of quadratic equations
and the odd degree of the field extension. All calculations below take
place in a fixed group over $\F_8$, independently of the ambient field.

\begin{theorem}\label{main}
Let $k=\F_{2^N}$ and $L=\F_{2^e}\subseteq k$, where $N$ is odd
and $3\mid e$. Suppose $a,b,c\in\Sz(k)$ have the following properties:
$a,c$ have odd order greater than one, $b$ has odd order greater than
one, and $ab,bc,abc$ are involutions. If
\[
 F=\langle a,b\rangle\cong\Sz(8),\qquad
 E=\langle b,c\rangle\cong\Sz(2^e)
\]
are genuine Suzuki subfield groups, then $F\le E$.
In particular they cannot satisfy $F\cap E=\langle b\rangle$.
\end{theorem}

Put $K=\F_8$ and $\theta=2^{(N+1)/2}$. The restriction of $\theta$
to $K$ is the fourth-power map. We use the standard matrices
$U(x,y)$, $z(r)=U(0,r)$, $H(\lambda)$ and $J$ from the preceding
notes. In particular the standard representation is symplectic for
$J$, and $c_2(g)=\Tr(g)^\theta$ for every Suzuki matrix $g$.

\begin{lemma}\label{descent}
Let $K\subseteq L\subseteq k$ be finite fields of characteristic two,
with $[k:K]$ odd.
Suppose $s,t\in k$ satisfy an affine equation
$\alpha s+\beta t+\gamma=0$ over $L$, where
$(\alpha,\beta)\ne(0,0)$, and a finite list of quadratic equations
$f_i(s,t)=0$ over $L$. Suppose their homogeneous parts of degree two
have coefficients in $K$ and have no common zero in $\mathbb P^1(K)$.
Then $s,t\in L$.
\end{lemma}
\begin{proof}
The homogeneous parts have no common projective zero over $L$ either.
Indeed a finite projective coordinate of a common zero is a root of
a nonzero polynomial of degree at most two over $K$. Its degree over
$K$ divides the odd number $[L:K]$, so it belongs to $K$.
The projective point at infinity is already a point over $K$.

Choose a point $(s_0,t_0)\in L^2$ on the affine line. Every point of
the line over $k$ has the form
\[
             (s,t)=(s_0,t_0)+T(\beta,\alpha).
\]
At least one $f_i$ has a nonzero quadratic leading coefficient after
this substitution, by the assertion just proved. Consequently $T$
satisfies a nonzero quadratic equation over $L$. Its degree divides
the odd number $[k:L]$, so $T\in L$ and the conclusion follows.
\end{proof}

\section*{Normalization and complete coverage}

Conjugate $F$ to its standard embedding over $K$. The possible
nontrivial odd orders in $F$ are $5,7,13$. Each occurs to the first
power in $|F|=29120$, so Sylow's theorem conjugates a subgroup of
any one of these orders to a chosen subgroup. The following explicit
automorphisms allow its generator to be chosen as well. Write
$K=\F_2[z]/(z^3+z+1)$ and set
\[
 b_5=z(1)J,\qquad b_7=H(z),\qquad b_{13}=z(1+z)J.
\]
These matrices have orders $5,7,13$, respectively. Direct multiplication
gives
\begin{align*}
 U(1,0)^{-1}b_5U(1,0)&=b_5^2,\\
 b_7^{(2)}=b_7^2,\qquad Jb_7J&=b_7^{-1},\\
 h^{-1}b_{13}^{(2)}h&=b_{13}^2,
       \qquad h=U(1,0)H(z)^3.
\end{align*}
Here $(2)$ means entrywise squaring. The powers of two generate
$(\mathbb Z/5)^\times$ and $(\mathbb Z/13)^\times$, while two and
minus one generate $(\mathbb Z/7)^\times$. Thus these automorphisms
cover every generator in each chosen subgroup. They extend to
automorphisms of the ambient Suzuki group and preserve the fields
$K$ and $L$. We may take $b$ to be one of the three displayed matrices.

Put $u=ab$, $v=bc$. As in the preceding note, all 455 involutions
of the standard $\Sz(8)$ are
\[
 z(r)\quad(r\ne0),\qquad
 U(x,y)^{-1}Jz(r)JU(x,y)\quad(x,y\in K,\ r\in K^\times). \tag{1}
\]
The characteristic-polynomial lemma and the odd order of $a$ give
$\Tr(bu)\ne0$. For $b_7$, discard also the fourteen involutions
which fix one of its two fixed ovoid points: otherwise $F$ would
fix that point. In coordinates these are the seven $z(r)$ and the
seven $Jz(r)J$. No other candidate is discarded.

Reduce the remaining list under conjugation by $b$, and also under
entrywise squaring when $b=b_5$, which has entries in $\F_2$.
The complete counts are as follows.
\[
\begin{array}{c|r|r|r}
 |b|&\text{eligible involutions}&\text{orbit size}&\text{orbits}\\\hline
 5&390&15&26\\
 7&378&7&54\\
 13&390&13&30
\end{array}                                                     \tag{2}
\]
The certificate checker regenerates (1), both exclusions, and every
orbit. The case files contain each representative and every element
of its orbit; their disjoint union is checked against the complete
eligible list. This supplies all 110 first-face cases without an
ambient-group enumeration.

\section*{The affine consequences}

Set $B=\Tr(b^2v)$, $Q=\Tr(bv)$, $B_1=B^\theta$ and
$Q_1=Q^\theta$. These elements belong to $L$ because $b,v\in E$.
The identity $Q=\Tr(c)$ gives $Q\ne0$. Write a general self-adjoint
matrix for the symplectic form $J$ as
\[
 V=\begin{pmatrix}
 v_0&v_1&v_2&v_3\\v_4&v_5&v_6&v_2\\
 v_7&v_8&v_5&v_1\\v_9&v_7&v_4&v_0
 \end{pmatrix}.
\]
The same eighteen linear equations as before are
\[
 ub^{-1}V=Vbu,\qquad \Tr(b^2V)=B,\qquad\Tr(bV)=Q.     \tag{3}
\]
In every case their rank in the ten matrix entries is eight.
Their complete solution is $V=V(B,Q,s,t)$, linear over $K$.
The two free entries $s,t$ are specified in each case file. The checker
verifies the entire linear system, its rank and these free entries.

Let $I_1,\ldots,I_m$ be the nonzero entries of $V^2-1$ followed by
$c_2(b^2V)-B_1$ and $c_2(bV)-Q_1$. They vanish at the hypothetical
configuration. Their terms of degree two in $s,t$ have coefficients
in $K$, independent of the trace values. For 106 cases the certificate
gives constants $c_i\in K$ and the identity
\[
 \sum_i c_iI_i=(a_0B+a_1Q)s+(d_0B+d_1Q)t+\gamma,
              \qquad \gamma\in K[B,Q,B_1,Q_1],         \tag{4}
\]
with $a_0d_1+a_1d_0\ne0$. The two linear coefficients in (4) cannot
vanish together because $Q\ne0$. For each case, the homogeneous
quadratic parts of the $I_i$ have no common zero in $\mathbb P^1(K)$.
This last statement is verified directly at its nine points
$[1:r]$, $r\in K$, and $[0:1]$; no conclusion over an extension is
inferred without Lemma~\ref{descent}.

All coefficients of the $I_i$ lie in $L$. Lemma~\ref{descent} applied
to (4) therefore gives $s,t\in L$ in these 106 cases, and every
entry of $v$ belongs to $L$.

\section*{The four exceptional cases}

The remaining cases have $|b|=13$. They are numbered 14, 17, 19 and
24 in the stored orbit list, with indices starting at zero. Each has
an affine consequence as in (4), but its coefficient matrix is
singular. Unless the two coefficients vanish together, the same
quadratic argument applies: the homogeneous forms again have no
common zero in $\mathbb P^1(K)$. The checker verifies this assertion
for these four cases as well.

We describe the branch in which both coefficients vanish. Encode
$a_0+a_1z+a_2z^2\in K$ by the integer $a_0+2a_1+4a_2$.
All arithmetic in the following two tables uses this encoding, not
integer addition or multiplication. The exact affine identities give
\[
\begin{array}{c|c|c|c}
 \text{cases}&B/Q&\text{remaining equation}&\text{possible }Q\\\hline
 14,17&4&4Q^2+7=0&7\\
 19,24&5&6Q^2+3Q^\theta+4=0&2,4
\end{array}                                                     \tag{5}
\]
For the first row, $Q^2=3=7^2$, so $Q=7$. For the second row,
$Q^\theta=2Q^2+5$. Applying $\theta$ and using
$(Q^\theta)^\theta=Q^2$ gives
$5Q^4+Q^2+7=0$, which is the square of
$3Q^2+Q+5=0$. Hence $Q$ has degree at most two over $K$.
Since $[k:K]$ is odd, $Q\in K$; substituting into the original
equation gives precisely $Q=2,4$. In all six resulting fibres,
$B,Q,B_1,Q_1$ thus belong to $K$.

After substituting these trace values, write the generators as
$\overline I_i\in K[s,t]$. The supplied identities include
$T(t)=\sum_i h_i\overline I_i$ and
$S(s,t)=\sum_i k_i\overline I_i$ with the following polynomials.
Here $T(t)$ is the product of the distinct linear factors
$t+r$ for $r$ in the indicated root set.
\[
\begin{array}{c|c|c|l}
 \text{case}&Q&\text{roots of }T&S(s,t)\\\hline
 14&7&1,5,6&s+2t^2+2t+6\\
 17&7&1,5,6&s+2t^2+2t+6\\
 19&2&6,7&s^2+4s+5t+3\\
 19&4&4&s+2\\
 24&2&2,3,6&s+5t^2+3t+1\\
 24&4&6&s
\end{array}                                                     \tag{6}
\]
The product formula forces $t\in K$. Each displayed $S$ is monic
of degree one or two in $s$, so the odd degree of $k/K$ gives
$s\in K$ as well. This completes the descent in every exceptional
branch.

We have proved $v\in\operatorname{Mat}_4(L)$ in all 110 cases.
The standard subfield intersection property gives
$v\in\Sz(L)$. Since $b\in\Sz(K)\le\Sz(L)$ and
$E=\langle b,v\rangle$ has order $|\Sz(L)|$, it follows that
$E=\Sz(L)$. Thus $F=\Sz(K)\le E$, proving Theorem~\ref{main}.

\section*{Verification and the remaining boundary}

For each of the 110 cases, the packet contains the complete linear
substitution, its full first-face orbit, every generator $I_i$, and
the exact constant combination giving the affine equation. The
quadratic coefficients and all nine projective directions are checked
anew. For each of the six exceptional fibres, every polynomial used
in (6) is accompanied by its full expression in the ideal generated
by the $\overline I_i$. The checker verifies those polynomial
identities and the displayed factorizations, rather than relying on
a computer algebra system's assertion about a Gr\"obner basis.
Its arithmetic uses the same retained standard-library polynomial
code as the previous packets. The field argument in
Lemma~\ref{descent} and the passage from matrix descent to subgroup
containment are written premises, not computational assertions.

Together with the two accepted small-face exclusions, this theorem
removes every cyclic-intersection configuration with a face $\Sz(8)$.
Indeed an order-four middle rotation forces coprime face degrees and
is excluded by the first note. An odd middle rotation has order
$5,7$ or $13$; if the other degree is coprime to three, the accepted
cyclic-overlap theorem leaves only order five, excluded by the second
note. Otherwise Theorem~\ref{main} applies. Exchanging the two faces
by $(a,b,c)\mapsto(c^{-1},b^{-1},a^{-1})$ covers the opposite order.

Both first-face degrees greater than three remain unresolved in the
F003 recovery. No whole-original, conjecture, Lean, or compression
credit is claimed. This author packet still requires distinct review.
\end{document}
